%% =============================================================================
%% sm_e4_mcmc.tex — SM-E, Subsection E.4: MCMC Configuration
%% Body content only (\subsection level); parent structure in sm_e.tex
%% =============================================================================

\subsection{MCMC configuration}
\label{sme:mcmc}

All models are fitted using Stan
\citep{CarpenterEtAl2017} with the No-U-Turn Sampler
\citep[NUTS;][]{Hoffman2014} via the \textsf{CmdStanR} interface
\citep{CmdStanR2024}.  \Cref{sme:tab-mcmc} reports the sampler
configuration.

\begin{table}[htbp]
  \centering
  \begin{tabular}{@{}l ll@{}}
    \toprule
    Setting & Empirical (M0--M3b) & Simulation \\
    \midrule
    Chains                  & 4 (parallel) & 4 (parallel) \\
    Warmup iterations       & 1{,}000 per chain & 1{,}500 per chain \\
    Sampling iterations     & 1{,}000 per chain & 2{,}000 per chain \\
    Total post-warmup draws & 4{,}000 & 8{,}000 \\
    \texttt{adapt\_delta}    & 0.95 & 0.95 \\
    \texttt{max\_treedepth}  & 12 & 14 \\
    \bottomrule
  \end{tabular}
  \caption{MCMC configuration.  The empirical case study (M0--M3b)
    uses the default \textsf{hurdlebb} settings; the simulation study
    uses longer runs and higher \texttt{max\_treedepth} to ensure
    adequate exploration across 200 replications.  In both cases,
    \texttt{adapt\_delta} $= 0.95$ reduces the step size to avoid
    divergent transitions in the state-varying coefficient models.}
  \label{sme:tab-mcmc}
\end{table}

\paragraph{Convergence targets.}
As stated in~\cref{alg:estimation}, we require the following
diagnostic thresholds for all parameters~\citep{Vehtari2021}:
(i)~split-$\hat{R} < 1.01$;
(ii)~minimum bulk $\ESS > 400$;
(iii)~minimum tail $\ESS > 400$; and
(iv)~zero divergent transitions.
All five models satisfy these thresholds; detailed diagnostics
are available in the replication package.

\paragraph{Warm-start initialization.}
The five-model sequence M0 $\to$ M1 $\to$ M2 $\to$ M3a $\to$ M3b
is fitted progressively: the posterior mean from model~$k$ is used
to initialize model~$k+1$.  Specifically:
\begin{itemize}
  \item Fixed-effect coefficients ($\balpha$, $\bbeta$,
    $\log\kappa$) are initialized at the predecessor's posterior
    means.
  \item Scale parameters ($\boldsymbol{\tau}$) are initialized at the
    predecessor's posterior means when dimensions match, or at~$0.1$
    when new dimensions are introduced.
  \item Cholesky factors ($\mathbf{L}_\varepsilon$) are initialized at
    the identity matrix, which corresponds to zero correlation---a
    conservative starting point that lets the sampler discover
    correlations during warmup.
  \item NCP variates ($\bz_s$) are initialized at zero, placing all
    states at the global mean.
  \item Policy moderator coefficients ($\bGamma$, M3b only) are
    initialized at the predecessor's posterior means or at zero for
    newly introduced parameters.
\end{itemize}
This strategy avoids the slow exploration that can occur when complex
models are initialized far from the posterior mode.

\paragraph{Runtime benchmarks.}
On a 16-core workstation, approximate wall-clock times are:
M0~($\sim$5 minutes),
M1~($\sim$15 minutes),
M2~($\sim$30 minutes),
M3a~($\sim$50 minutes),
M3b~($\sim$60 minutes),
and M3b-W~($\sim$70 minutes, including score extraction in the
\texttt{generated quantities} block).  The total sequential fitting
time is approximately 4~hours with warm-start initialization, compared
to approximately 7~hours with cold starts.
